\section{Related Work}

\paragraph{Chain-of-Thought Prompting.}
Recently,~\cite{wei2022CoT} has proposed chain-of-thought (CoT) prompting to improve the reasoning ability of LLMs. Specifically, CoT prompts LLMs to generate a series of intermediate steps that lead to the final answer of a multi-step problem. Most earlier work primarily concentrates on two main aspects: prompt design and decoding strategies. Zero-shot CoT~\citep{kojimalarge} employs the trigger sentence ``Let’s think step by step'' to provide guidance for the decoding of LLMs. Advanced sampling strategies have been explored to improve CoT by generating diverse reasoning paths, e.g., Self-Consistency~\citep{wang2022self}, Auto-CoT~\citep{zhang2022automatic}, Active-Prompting~\citep{diao2023active}, Complexity-based Consistency~\citep{fu2022complexity}, Multi-Chain Reasoning~\citep{yoran2023answering}, and Progressive-Hint Prompting~\citep{zheng2023progressive}. 

With the emergence of powerful LLMs, approaches based on self-evaluation have attracted increasing attention. These approaches involve the generation of initial output, followed by evaluating the output to acquire feedback, which is then utilized to refine the output. Evaluation feedback can come from the model itself, e.g., Self-refine~\citep{madaan2023selfrefine} and Tree of Thoughts~\citep{yao2023tree}) or external environments, e.g., QAaP~\citep{zhu2023question} and Reflection~\citep{shinn2023reflexion}. The intuition behind these approaches involves the utilization of robust LLMs to mimic the human cognition process. 

\paragraph{Generative Agents.}
Recently, LLM-based multi-agent intelligent, e.g., Generative Agents~\citep{park2023generative}, Ghost in the Minecraft~\citep{zhu2023ghost}, GPT-Bargaining~\citep{fu2023improving}, has drawn significant attention for enabling simulations of human behavior. Our work follows this research line to address the DoT problem of LLMs. Concurrent with our work, a few studies~\citep{xiong2023diving,du2023improving} also explore the multi-agent debate framework to  enhance the reasoning ability of LLMs. The main differences between our MAD framework and these works are: 
(1) we introduce an additional judge with an adaptive break mechanism to decide the optimal moment to conclude the debate;
(2) our work aims to address the DoT problem, which is an inherent deficiency of LLMs; 
and (3) we empirically find that our MAD framework can yield enhanced performance by employing agents with the identical backbone LLM.
